%% ===================================================================
\section{Antecedentes}\label{sec:antecedentes}

\subsection{Perturbación proactiva y contramedidas}

En esta revisión, una \emph{perturbación proactiva} es una modificación que el titular aplica a sus datos visuales antes de publicarlos, con el fin de degradar la utilidad de un modelo discriminativo entrenado o ajustado con ellos sin autorización. La literatura la denomina también ejemplo no aprendible, dato no explotable o ataque de disponibilidad. El concepto excluye dos fenómenos próximos: las perturbaciones que actúan solo durante la inferencia de un modelo ya entrenado y las protecciones dirigidas exclusivamente contra modelos generativos.

Una \emph{contramedida} es cualquier procedimiento que el adversario aplica a los datos protegidos o a su propio entrenamiento para recuperar el rendimiento del modelo. Incluye las transformaciones y compresiones de la entrada, los aumentos de datos, el entrenamiento adversario, la purificación mediante modelos generativos, los cambios del paradigma de aprendizaje y la detección y descarte de las muestras protegidas. La \emph{recuperación autorizada} es un caso distinto: un actor legítimo dispone de información o recursos del titular para restituir la aprendibilidad del dato.

\subsection{Condiciones de comparación y medidas de recuperación}\label{sec:condiciones}

La evaluación de una contramedida requiere comparar el rendimiento de modelos entrenados bajo cuatro condiciones:
\begin{itemize}
\item con datos limpios bajo el régimen estándar;
\item con datos protegidos sin la contramedida;
\item con datos protegidos después de aplicar la contramedida;
\item con datos limpios sometidos a la misma contramedida.
\end{itemize}
A partir de ellas se definen tres magnitudes, orientadas según la dirección de la métrica:
\begin{itemize}
\item la \emph{recuperación}, diferencia entre el rendimiento con contramedida y el rendimiento protegido sin ella;
\item la \emph{brecha al entrenamiento limpio}, diferencia entre el rendimiento limpio y el obtenido con contramedida;
\item la \emph{brecha al control con contramedida}, diferencia respecto del entrenamiento limpio sometido a la misma contramedida.
\end{itemize}
Cuando la métrica es la exactitud en porcentaje, las tres se expresan en puntos porcentuales (pp).

Las magnitudes se interpretan en conjunto:
\begin{itemize}
\item Una recuperación alta con una brecha baja al entrenamiento limpio indica que la contramedida neutraliza la protección.
\item Una brecha alta al entrenamiento limpio con una brecha baja al control indica que la pérdida remanente se debe a la propia contramedida y no a la protección.
\end{itemize}

\subsection{Revisiones relacionadas}

Tres trabajos de síntesis abordan el tema desde perspectivas distintas (Tabla~\ref{tab:revisiones}). Li et al.~\cite{li2025survey} revisan los métodos de generación de datos no aprendibles, sus benchmarks, métricas y aplicaciones, y comparan los enfoques en términos de no aprendibilidad, imperceptibilidad, eficiencia y robustez. Zhang et al.~\cite{zhang2026sok} relacionan la no aprendibilidad con el desaprendizaje de máquinas y aportan una evaluación empírica propia. Cinà et al.~\cite{cina2023wild} sistematizan los ataques y defensas de envenenamiento del entrenamiento, dentro de los cuales la perturbación proactiva aparece como un caso de envenenamiento indiscriminado. Los dos primeros trabajos son preimpresiones.

Ninguno de estos trabajos sigue un protocolo sistemático con criterios de calidad declarados ni cuantifica la recuperación del adversario frente a las protecciones. Esa síntesis cuantitativa, restringida a las comparaciones que pueden interpretarse, es el aporte de la presente revisión.

\begin{table}[t]
\caption{Revisiones relacionadas y diferencias con la presente revisión}\label{tab:revisiones}
\footnotesize
\begin{tabularx}{\textwidth}{@{}L{3.0cm}Y{1.0}Y{1.0}@{}}
\toprule
Revisión & Alcance & Diferencia con esta revisión \\
\midrule
Li et al. (2025) \cite{li2025survey} & Métodos de generación de datos no aprendibles, benchmarks, métricas, fundamentos teóricos y aplicaciones; compara enfoques por no aprendibilidad, imperceptibilidad, eficiencia y robustez & Revisión narrativa, sin protocolo sistemático ni evaluación de calidad; no cuantifica la recuperación del adversario \\
Zhang et al. (2026) \cite{zhang2026sok} & Sistematización conjunta de no aprendibilidad y desaprendizaje de máquinas, con evaluación empírica propia y una garantía teórica sobre la desmemorización & Analiza la relación entre ambas defensas; no sintetiza la evidencia publicada sobre contramedidas \\
Cinà et al. (2022) \cite{cina2023wild} & Ataques y defensas de envenenamiento del entrenamiento en aprendizaje automático, con énfasis en visión por computadora & Trata la perturbación proactiva como un caso del envenenamiento indiscriminado; no aborda su uso como protección de datos \\
\botrule
\end{tabularx}
\end{table}

